\begin{abstract}
Language models can be wired together through their hidden states, and several systems report that such latent channels help. Two questions are rarely separated: how much of a message's benefit is content the receiver lacks, and does building the channel require real text from the models' training distribution? We study both with an information-asymmetric protocol: a sender reads a prefix the receiver never sees, the receiver gets the sender's hidden states through a learned bridge, and donor messages from other windows decompose the receiver's gain into generic, article-level and window-specific parts. The choice of donor is decisive: swapping messages within a batch, where most donors come from the same document, attributes 54\% of a bridge's benefit to generic context; with donors from other documents the generic part is 6\%. Between independently pretrained Pythia models, a linear bridge fitted on real text gives the receiver about the same content gain as the receiver's own prefix states injected the same way (ratio 0.99 [0.94, 1.03]). A bridge fitted only on encyclopedia-style text written by an LLM with no corpus access, with layer selection also done on such text, reaches 82\% [79, 85] of that reference's content gain; bridges fitted on the models' own samples reach 44--59\%, and a weight-only least-squares map 2\%. In the main WikiText experiment, bridges fitted with the tested synthetic-text recipes leave the receiver worse off on average than no message when the message comes from a different document, unlike the real-text bridge. \SCALESENTENCE
\end{abstract}
\begin{center}\small Code, results and fitted bridges: \url{https://github.com/dongzhaohe321418-lab/content-or-context}\end{center}

\section{Introduction}
\label{sec:intro}

Several recent systems let language models communicate through hidden states rather than text: a sender's key--value cache is fused into a receiver's \citep{fu2026c2c}, agents built on one backbone share last-layer states and caches \citep{zou2025latentmas}, a receiver reads a sender's last-layer states through a trained adapter \citep{du2025interlat}, and affine maps transfer features between residual streams \citep{chen2025stitching}. Two questions are often left implicit. First, \emph{how much of the benefit is content?} A receiver handed a plausible representation of some preceding text may improve simply because its input now looks like a document in progress, whatever that document is. Second, \emph{what does building the channel require?} The projectors in these works are trained on instruction, task or pretraining-style data, or (in LatentMAS) restricted to agents that share a backbone. Connecting models whose data cannot be shared requires building the bridge from something else.

We study both questions in a setting where they are well posed: causal language models pretrained in separate runs (Pythia 70M, 160M and 410M, which share a tokenizer and training corpus; \citealp{biderman2023pythia}), \XFAMILYINTRO\ The sender $A$ reads a 128-token prefix; the receiver $B$ reads only the next 32 tokens and receives $A$'s prefix states, mapped by a bridge and injected into its residual stream (\S\ref{sec:protocol}). We measure the reduction in $B$'s cross-entropy on the suffix and decompose it with \emph{donor messages}: the same message source applied to a prefix from another article (generic part), or from another window of the same article (separating article-level from window-specific content). The \emph{receiver-state reference}, $B$'s own prefix states injected the same way, serves as a yardstick, and $B$ reading the prefix as text as a second one.

Bridges are fitted by matching $B$'s own predictive distribution, which needs probe text but no labels. The conditions differ in the probe text (and, unavoidably, in how much of it there is): real in-domain text; encyclopedia-style articles written by an LLM with no tools and no corpus access; unconditional samples from the two models; or none at all (a least-squares map between the models' embedding matrices). The LLM condition is motivated by recent evidence that frontier models can reason about network weights directly \citep{takehi2026sparks}; in small testbeds an LLM could design a complete bridge from weights alone (Appendix~\ref{app:small}). Here we ask how far an LLM's world knowledge goes as a substitute for data at the scale of pretrained models.

\paragraph{Contributions.}
\begin{itemize}
\item \textbf{A control that separates content from context, and a pitfall.} In the main pair, donor messages from other documents show that almost all of a latent message's benefit is content (generic part 2--7\% of the benefit of text, the receiver-state reference and real-text bridges). Donors from the same document, as produced by swapping messages within sorted batches, remove about half of the content and misattribute it to generic context (\S\ref{sec:generic}).
\item \textbf{Bridges fitted on synthetic text work, with a measurable cost.} An LLM-written-text bridge, built and selected without corpus text, reaches 82\% of the receiver-state reference's content gain, against 44--59\% for self-samples and 2\% for weight algebra. In the main WikiText experiment, unlike the real-text bridge, the tested synthetic-text bridges have a negative average generic part: a message about a different document leaves the receiver worse off on average than no message (\S\ref{sec:datafree}). \SCALECONTRIB
\item \textbf{Where and how latent messages fall short.} Injecting deeper sharply reduces what even the receiver's own prefix states deliver; at the shallowest depth they deliver less content than text, and for every tested number of transmitted positions the bridged message is worth less to this receiver than the same number of text tokens (\S\ref{sec:where}, \S\ref{sec:bandwidth}). Hand-over stitching, which replaces the receiver's lower computation, measures something different and should not be read as a test of communication (\S\ref{sec:diag}).
\end{itemize}
